%!TEX root = ../../../main.tex

\chapter{Diskussion}
\label{chap:diskussion}

\paragraph{Code entscheidet, das Modell liest.} Über alle Versionen stimmte fast immer, was Code entscheidet: Fristen, Listenfilter, Bestand. Fehler entstehen, wo das Modell lesen muss: Es nimmt die falsche Spalte, verwechselt ähnliche Katalogeinträge oder beschreibt einen Filter falsch. Deshalb habe ich dem Modell schrittweise Aufgaben abgenommen. P2, der alles selbst entscheidet, liegt deutlich darunter. Für ein 7--8B-Modell auf der CPU ist ein fester Ablauf mit Prüfung besser als freie Werkzeugwahl.

\paragraph{Probleme während der Entwicklung.} \emph{Wortlisten:} Meine erste Bestandsprüfung kannte nur die Dokumentarten, an die ich gedacht hatte. Heute leitet der Katalog den Wortschatz aus den Dokumenten ab. Aber \enquote{Haftpflicht} steht im Kfz-Versicherungsschein, eine private Haftpflichtversicherung gibt es trotzdem nicht. Jetzt prüft Code auch das Eigenschaftswort vor der Dokumentart. \emph{Latenz:} In einer Messung dauerte jeder Plan 100\,s statt 17\,s. Ein zweiter Prompt-Cache-Platz half nicht, denn die Ursache war der Messaufbau: Der Lauf wechselte je Frage zwischen P1 und P2, und das Modell von P2 verdrängte das von P1 aus dem Speicher. \emph{Prüfungen, die zufällig passen:} Wie T14 zeigt (Abschnitt~\ref{sec:weg}), kann eine Zahlprüfung aus dem falschen Grund bestehen. Umgekehrt meldete meine Bewertung bei \enquote{1. Januar 2024} eine erfundene Zahl, weil sie nur die Schreibweise 01.01.2024 kannte. Beides fiel erst beim Lesen einzelner Antworten auf, nicht an der Kennzahl.

\paragraph{Grenzen.} Alle Dokumente und Testfragen stammen von mir, die Sätze sind klein, und gemessen habe ich nur auf einem Laptop.
